\documentclass[11pt,a4paper]{article}

\usepackage[T1]{fontenc}
\usepackage[utf8]{inputenc}
\usepackage{lmodern}
\usepackage[a4paper,margin=26mm]{geometry}
\usepackage{amsmath,amssymb,amsthm,mathtools}
\usepackage{booktabs,array}
\usepackage{microtype}
\usepackage{xcolor}
\usepackage{enumitem}
\usepackage{url}
\usepackage[colorlinks=true,linkcolor=blue!55!black,citecolor=blue!55!black,
            urlcolor=blue!55!black]{hyperref}

\newtheorem{theorem}{Theorem}
\newtheorem{lemma}[theorem]{Lemma}
\newtheorem{proposition}[theorem]{Proposition}
\newtheorem{remark}[theorem]{Remark}
\newcommand{\bits}{\{0,1\}}
\newcommand{\concat}{\mathbin\Vert}
\newcommand{\sample}{\mathrel{\leftarrow\$}}
\newcommand{\KeyGen}{\ensuremath{\mathsf{KeyGen}}}
\newcommand{\Sign}{\ensuremath{\mathsf{Sign}}}
\newcommand{\Verify}{\ensuremath{\mathsf{Verify}}}
\newcommand{\PRF}{\ensuremath{\mathsf{PRF}}}
\newcommand{\PRFmsg}{\ensuremath{\mathsf{PRF}_{\rm msg}}}
\newcommand{\Hmsg}{\ensuremath{\mathsf{H}_{\rm msg}}}
\newcommand{\FORC}{\ensuremath{\mathsf{FORC}}}
\newcommand{\HT}{\ensuremath{\mathsf{HT}}}
\newcommand{\pk}{\ensuremath{\mathsf{PK}}}
\newcommand{\sk}{\ensuremath{\mathsf{SK}}}
\newcommand{\Prb}{\mathop{\rm Pr}}
\newcommand{\E}{\mathbb E}

\title{\bfseries Concrete FORC Lifetime Failure in CEDRUS-$\alpha$:\\
Ordinary Fresh-Message Forgery within the NGCC\\
Chosen-Message Query Budget}
\author{Martin Feussner\\
\small Selmer Center, University of Bergen\\
\small \texttt{martin.feussner@uib.no}}
\date{3 October 2026}

\begin{document}
\maketitle

\begin{center}
\fbox{\parbox{0.92\linewidth}{\small\textbf{Provenance and review status.}
The candidate-specific parameter analysis, exact concentration argument,
implementations, experiments, release package, and initial manuscript were
produced by OpenAI Codex using Daybreak Blue at Ultra reasoning effort during
an AI-run audit of NGCC signature candidates.  A separately instructed
hostile-review agent reconstructed the specification and attack independently,
recomputed the bounds, implemented a second reduced verifier, and accepted the
result only with the scope and resource qualifications stated below.  The
generic FORC accumulation mechanism is already described by the CEDRUS-$\alpha$
submission, and Abri--Katz already account for public hash trials; no novelty
is claimed for either idea.  At the date above, no human cryptanalyst has
independently reviewed the proof, implementations, measurements, or novelty
assessment.  This report is a preliminary disclosure for human review and
reproduction.}}
\end{center}
\vspace{0.7em}

\begin{abstract}
CEDRUS-$\alpha$ is a stateless hash-based signature whose bottom few-time
scheme, FORC, reveals intermediate nodes of forward hash chains.  Its own
Section~3.2 recognizes that signatures which reuse one bottom FORC address can
accumulate enough openings to cover a new digest.  We evaluate the consequence
after adding the public target-hash trials that an ordinary chosen-message
forger may perform.

The adversary collects independently randomized signatures and groups them by
the complete $h$-bit bottom address.  For each address, coordinate, and leaf it
retains the earliest exposed chain node and that node's authentication path.
For a message never queried to the signer, it enumerates serialized public
$R$ values until $\Hmsg(R,\pk.\mathsf{seed},\pk.\mathsf{root},M)$ selects only
covered coordinates at one retained address.  Each node is advanced publicly,
the coordinate openings are mixed, and one hypertree suffix for that address
authenticates the resulting fixed FORC public key.  This is an ordinary
fresh-message forgery: it requires no reset, repeated randomness, deterministic
mode, fault, related key, weak key, or implementation defect.

We derive exact first and second coverage moments.  Poissonizing acquisition
makes address buckets independent, so a Chebyshev bound controls the realized
table rather than merely its expectation.  All eight submitted parameter sets
admit constant-success strategies with signing means at most $2^{77.520}$ and
attacker public-$\Hmsg$ work $Q+T$ below $2^{77.56}$, with a hard abort before
$2^{80}$ signing requests.  Here the attacker performs one $\Hmsg$
recomputation for each acquired signature plus $T$ target trials.  For 160f,
retaining one fixed six-bit address
prefix gives signing mean $2^{71.458}$, fixed public-hash budget $2^{67.160}$,
attacker $\Hmsg$ work $Q+T=2^{71.530}$, expected full-path table size $2^{77.479}$
bytes, and deterministic table capacity $2^{79.175}$ bytes.  The physical
qualification is substantial: expected signature traffic is $2^{85.703}$
bytes and the honest signing service performs about $2^{89.060}$ hash
evaluations.  Across all eight $Q+T$-optimized rows, the largest expected
resource exponent in its own natural unit is $100.852$, for honest-signer hash
calls in 512s.  These heterogeneous resources are not combined into one cost.

Two independent reduced implementations produced accepted fresh-message
forgeries and rejected wrong-message, wrong-$R$, wrong-address, backward-chain,
authentication-path, suffix, and whole-transcript controls.  A repaired native
256f build separately confirms the normative address, PRF, WOTS, signing, and
verification behavior.  An all-eight-set full-parameter FORC-level experiment
also confirms coordinate splicing under both endpoint conventions.  It is not
a full-parameter signature forgery, which was not run.  The narrow new
conclusion is a CEDRUS-$\alpha$
parameter and missing-lifetime-policy failure under the official NGCC query
budget, not a new generic accumulation or public-grinding technique.
\end{abstract}

\section{Introduction}

CEDRUS-$\alpha$ is a stateless hash-based signature in the SPHINCS family
\cite{cedrus-spec,cedrus-framework}.  A hypertree authenticates the public key
of a bottom few-time signature.  CEDRUS-$\alpha$ replaces the usual FORS
bottom layer with FORC, whose leaves are endpoints of short forward hash
chains.  A FORC signature reveals an intermediate node selected by the message
digest.  Later positions on the same chain become computable, so the bottom key
degrades gradually as its address is reused.

This degradation is not hidden in the submission.  Section~3.2 explicitly
mixes over the number of signatures reaching one FORC instance and gives an
approximate probability that those openings cover one new digest.  Its
displayed quantity has no public-hash-query variable.  A forger can hold a
fresh message fixed and enumerate the serialized randomizer $R$, which is a
public input to the verifier's message digest.  Algorithm~21 reads $R$ from the
signature and recomputes the digest; it has neither $\sk.\mathsf{prf}$ nor
\textsf{addrnd} and makes no $\PRFmsg$-origin check.  The one-target probability
is therefore amplified over
$q_H$ distinct public inputs.  Abri and Katz give exactly this style of
$q_H$ accounting for a related hash-based construction
\cite[Eq.~(2)]{abri-katz}; we do not claim it as a new technique.

The concrete consequence for the submitted parameters was left unstated.  The
submission gives no maximum number of signatures per key, while the official
NGCC evaluation model allows up to $2^{80}$ chosen-message signing queries
\cite{ngcc-criteria}.  Once acquisition and target grinding are optimized
together, every submitted parameter set admits a fresh-message forgery below
that query ceiling.

Our contributions are as follows.
\begin{enumerate}[leftmargin=*,itemsep=2pt]
  \item We give an explicit coordinate-mixing forgery against the coherent
        specification and prove acceptance by the ordinary verifier.
  \item We derive exact first and second moments of the target coverage.  A
        Poissonized acquisition algorithm, a hard $2^{80}$ abort, and a
        concentration bound yield a fixed target-search budget and a rigorous
        constant success probability.
  \item We evaluate all eight submitted parameter sets and keep signing
        queries, public-hash trials, response traffic, attacker storage and
        processing, and honest signing-service work as separate resources.
  \item We give a six-bit address-prefix shard whose fully provisioned simple
        donor table is below $2^{80}$ bytes, plus an independent fixed-address
        low-memory analysis.
  \item We reproduce the complete attack at reduced dimensions using ordinary
        hedged signatures and extensive negative controls.  We separately
        validate a repaired native 256f core and an all-eight-set,
        full-parameter FORC splice, excluding the known submitted source
        deviations while avoiding any claim of a full-size forgery run.
  \item We delimit novelty to the exact CEDRUS-$\alpha$ parameter/lifetime
        consequence.  FORC accumulation and the $q_H$ multiplier are prior
        art.
\end{enumerate}

The result is a specification-level parameter and key-lifetime failure in the
ideal-hash model.  It is not machine-practical: the signing-query and
public-$\Hmsg$ counts sit beside very large responses and very expensive honest
signing.  We report those resources separately and make no claim that the
full-parameter attack was executed on the audit machine.

\section{Relevant construction}

\subsection{Keys, randomized signing, and digest parsing}

For node length $n$, key generation produces
\begin{equation}
 \sk=(\sk.\mathsf{seed},\sk.\mathsf{prf},
      \pk.\mathsf{seed},\pk.\mathsf{root}),\qquad
 \pk=(\pk.\mathsf{seed},\pk.\mathsf{root}).             \label{eq:keys}
\end{equation}
The first three components are independently sampled $n$-byte strings, and
$\pk.\mathsf{root}$ is the hypertree root.

Ordinary hedged signing samples a fresh $n$-byte \textsf{addrnd} and computes
\begin{align}
 R&=\PRFmsg(\sk.\mathsf{prf},\mathsf{addrnd},M),
                                                               \label{eq:R}\\
 \mathsf{digest}&=\Hmsg(R,\pk.\mathsf{seed},
                         \pk.\mathsf{root},M).                 \label{eq:hmsg}
\end{align}
The digest parses as $k$ FORC pairs
$(x_i,\ell_i)\in[0,2^a)\times[0,W)$, a bottom-tree index
$\mathsf{idxTree}$ of $h-h_0$ bits, and a bottom-leaf index
$\mathsf{idxLeaf}$ of $h_0$ bits, where $W=w'$.  We write
\begin{equation}
             A=(\mathsf{idxTree},\mathsf{idxLeaf})             \label{eq:A}
\end{equation}
for the complete $h$-bit bottom FORC address.  The verifier receives $R$ in the
signature: Algorithm~21 applies \textsf{SIG.getR()} and
recomputes~\eqref{eq:hmsg}.  Forgery validity is determined by this public
verification algorithm.  It has no secret input with which to test
relation~\eqref{eq:R}, so any serialized $n$-byte $R^*$ is syntactically
eligible and the attacker may grind distinct $R^*$ values on the fresh target
message.

Table~\ref{tab:params} gives the parameters used below.  The final column is
the specification's approximate number of hash evaluations for one honest
signature, used only to expose signing-service cost.

\begin{table}[ht]
\centering
\small
\caption{Submitted CEDRUS-$\alpha$ parameters.}
\label{tab:params}
\begin{tabular}{lrrrrrrrrr}
\toprule
set & $n$ & $h$ & $d$ & $a$ & $k$ & $W$ & OTS len & $|\sigma|$ & sign hashes\\
 & bytes &&&&&&& bytes & \\
\midrule
160s &20&68&9 &10&16&8&30 &10,300 &2,743,319\\
160f &20&66&17&7 &28&4&40 &19,420 &199,011\\
256s &32&67&9 &12&23&4&48 &25,568 &4,163,230\\
256f &32&65&14&8 &45&2&63 &43,296 &456,199\\
384s &48&65&8 &12&38&4&88 &60,672 &5,347,900\\
384f &48&65&12&10&51&2&80 &76,176 &1,490,831\\
512s &64&65&8 &13&47&4&101&98,048 &10,555,995\\
512f &64&66&11&10&66&4&109&127,488&2,481,272\\
\bottomrule
\end{tabular}
\end{table}

\subsection{FORC equations}

Put $B=2^a$.  For coordinate $i$, leaf $x$, and chain position $j$, let the
global chain index be $u=iB+x$.  At complete address $A$, the coherent
Algorithms~13--18 define
\begin{align}
 s_{i,x}&=\PRF(\pk.\mathsf{seed},\sk.\mathsf{seed},
                  \mathsf{ADRS}_{\rm FTS\mbox{-}PRF}(A,u)),     \label{eq:secret}\\
 C_{i,x}(0)&=s_{i,x},\nonumber\\
 C_{i,x}(j+1)&=F(\pk.\mathsf{seed},
                  \mathsf{ADRS}_{\rm FTS\mbox{-}CHAIN}(A,u,j+1),
                  C_{i,x}(j)),                                 \label{eq:chain}\\
 L_{i,x}&=C_{i,x}(W).                                          \label{eq:leaf}
\end{align}
The $B$ leaves $L_{i,x}$ form a height-$a$ Merkle tree with root $T_i$.
The FORC public key is
\begin{equation}
 \mathsf{PK}_{\FORC}(A)=T_k(\pk.\mathsf{seed},
      \mathsf{ADRS}_{\rm FTS\mbox{-}ROOT}(A),
      T_0\concat\cdots\concat T_{k-1}).                       \label{eq:forcpk}
\end{equation}

For digest coordinate $(x_i,\ell_i)$, the signature contains
$C_{i,x_i}(\ell_i)$ and the $a$ Merkle siblings for leaf $x_i$.
Verification applies $W-\ell_i$ public forward steps, reconstructs $T_i$ with
that path, and evaluates~\eqref{eq:forcpk}.  A hypertree suffix signs
$\mathsf{PK}_{\FORC}(A)$ at the same complete address and authenticates it to
$\pk.\mathsf{root}$.

\subsection{Resolution of specification ambiguities}
\label{sec:ambiguities}

The PDF contains stale fields and names, but no ambiguity supplies the attack.
We use the following coherent reading.
\begin{enumerate}[leftmargin=*,itemsep=2pt]
  \item Equation~(1.3), the $\Hmsg$ rows of Tables~1.2--1.3, and
        Figure~1.14 retain a dangling public-counter formalism.  The exact token
        \textsf{ctrFTS} occurs only in Algorithm~20, line~16, where an undefined
        value is appended.  Algorithm~20, line~4 computes $\Hmsg$ without it;
        Algorithm~21 neither parses nor checks it, computes $\Hmsg$ without it,
        and gives a counter-free length equation.  Table~1.1 and
        Equation~(4.1) give the same counter-free sizes.  No generation loop,
        acceptance predicate, rejection rule, or counter distribution is
        specified.  We therefore follow the executable, counter-free
        Algorithms~20--21 and published sizes.  An unconstrained public counter
        would add another public grinding input and four response bytes; it
        would not defeat the attack.
  \item Stale \textsf{FORS} and \textsf{PORS} names in Algorithms~20--21 mean
        the FORC algorithms defined in Section~1.9.
  \item The preliminary digest-length formula omits the $k\log_2 W$ position
        bits.  The later formula, parser, and Algorithms~16--21 include them.
  \item Algorithms~14--18 operationally disclose positions
        $0,\ldots,W-1$ and authenticate $C(W)$, while the prose can be read as
        authenticating $C(W-1)$.  Write $E\in\{W-1,W\}$ for the chosen endpoint.
        In either reading, a donor can reach a target exactly when
        $\ell_{\rm donor}\leq\ell_{\rm target}$: it advances by their difference
        and verification applies the remaining $E-\ell_{\rm target}$ steps.
        Relabelling positions by one also preserves this order.  Consequently
        $g_r$, the pairwise second moment, and every concrete table below are
        invariant under the endpoint convention.
\end{enumerate}
Reading the stale counter as evidence for an unpublished forced-pruning rule
would define a different scheme.  Ordinary ideal-$\Hmsg$ chunks are uniform
and independent in the submitted algorithms.  All sixteen reference and
optimized source trees corroborate the counter-free reading: signing serializes
only $R$ before the FORC bytes, verification consumes that $R$, the message hash
input is $R\concat\pk\concat M$, and \textsf{SPX\_BYTES} contains no four-byte
counter term.

\section{The coordinate-mixing forgery}

\begin{lemma}[Forward completion]
Fix a key, complete address $A$, coordinate $i$, and leaf $x$.  An opening at
position $\ell$ constructs a verifier-valid opening for every target position
$\ell'\geq\ell$ at the same $(A,i,x)$.
\end{lemma}
\begin{proof}
Apply~\eqref{eq:chain} exactly $\ell'-\ell$ times to obtain
$C_{i,x}(\ell')$.  Keep the donor's authentication path.  The verifier applies
the remaining $E-\ell'$ steps and reaches the same fixed endpoint $C_{i,x}(E)$,
where $E=W$ under Algorithms~14--18 and $E=W-1$ under the shifted prose
reading.  It therefore reconstructs the same leaf and root $T_i$ in either
convention.
\end{proof}

\begin{lemma}[Coordinate mixing]
At a fixed complete address $A$, valid openings for different coordinates may
be taken from different signatures and combined with one donor hypertree
suffix.
\end{lemma}
\begin{proof}
For a fixed key and $A$, each $T_i$ in~\eqref{eq:forcpk} is fixed independently
of the digest that selected an opening.  Every valid coordinate opening
reconstructs its fixed $T_i$, so their ordered compression is the fixed
$\mathsf{PK}_{\FORC}(A)$.  Every honest suffix at $A$ authenticates exactly
that value.  The proof never transplants a path across leaves, coordinates, or
addresses.
\end{proof}

The adversary runs the following algorithm.
\begin{enumerate}[leftmargin=*,itemsep=2pt]
  \item Query ordinary randomized signatures on distinct chosen messages.
        Recompute their public digests and group them by the complete address
        $A$.
  \item For every observed $(A,i,x)$, retain the donor with the smallest
        $\ell$, its node, and its complete authentication path.  Retain one
        hypertree suffix for each occupied address.
  \item Choose a message $M^*$ never sent to the signing oracle.  Enumerate
        distinct $n$-byte strings $R^*$ and evaluate
        $\Hmsg(R^*,\pk.\mathsf{seed},\pk.\mathsf{root},M^*)$.
  \item Stop when its address $A^*$ is retained and every selected
        $(x_i^*,\ell_i^*)$ has a donor at $(A^*,i,x_i^*)$ with
        $\ell_i\leq\ell_i^*$.
  \item Advance every donor to its target position, copy its path, and attach
        one retained suffix for $A^*$.
\end{enumerate}

\begin{theorem}[Fresh-message acceptance]
The output of the algorithm is accepted by the specification-conformant
verifier on $M^*$.
\end{theorem}
\begin{proof}
The verifier recomputes the target digest from the serialized $R^*$ and fresh
$M^*$.  The first lemma shows that each constructed coordinate reaches the
correct fixed endpoint and root.  The second lemma shows that their compression
is the fixed $\mathsf{PK}_{\FORC}(A^*)$ and that the copied suffix authenticates
it to the public root.  All verifier checks therefore pass.  Freshness follows
because $M^*$ was never queried; $R^*$ need not satisfy the signer's secret
$\PRFmsg$ relation.
\end{proof}

Acquisition uses independent fresh signer randomness.  The attack does not
ask for, cause, or test an $R$ collision.  Complete messages and complete FORC
vectors need not repeat.

\section{Exact coverage and concentration}

\subsection{One address}

Suppose $r$ signatures reach one address.  For one target coordinate with
position $y$, a uniform donor covers it when it chooses the same leaf and one
of positions $0,\ldots,y$, with probability $(y+1)/(BW)$.  Averaging over a
uniform target position gives the exact mean covered fraction
\begin{equation}
 g_r=\frac1W\sum_{y=0}^{W-1}
        \left[1-\left(1-\frac{y+1}{BW}\right)^r\right].       \label{eq:gr}
\end{equation}
The $k$ coordinate streams are independent conditional on $r$, so the mean
covered full-vector fraction is $g_r^k$.

Concentration across addresses needs a second moment.  Let $C_r$ be the
covered fraction of one coordinate's $BW$-element target domain, and define
$m_r(s)=(1-s/(BW))^r$.  For two ordered targets at positions $y,z$, the
probability that both donor sets have been hit is
\begin{equation}
 \begin{cases}
  1-m_r(y+1)-m_r(z+1)+m_r(\max\{y+1,z+1\}),&x=x',\\
  1-m_r(y+1)-m_r(z+1)+m_r(y+z+2),&x\ne x'.
 \end{cases}                                                \label{eq:pair}
\end{equation}
Averaging~\eqref{eq:pair} over the $B^2W^2$ ordered pairs gives
$s_r=\E[C_r^2]$.  Independence across coordinates gives second moment $s_r^k$
for the full-vector covered fraction.

\subsection{All addresses}

Let the adversary first sample its acquisition count
$L\leftarrow\mathsf{Pois}(Q)$ and abort if $L>2^{80}$.  Under the ideal digest
model, Poisson splitting makes the $2^h$ complete-address loads independent
$\mathsf{Pois}(\lambda)$ variables, where $\lambda=Q/2^h$.  Define
\begin{equation}
 \mu=\E_{S\leftarrow\mathsf{Pois}(\lambda)}[g_S^k],\qquad
 \nu=\E_{S\leftarrow\mathsf{Pois}(\lambda)}[s_S^k].          \label{eq:munu}
\end{equation}

For realized tables $Z_A$, one public target succeeds with probability
\begin{equation}
                  P=2^{-h}\sum_A Z_A.                        \label{eq:P}
\end{equation}
Independence of the buckets and $0\leq Z_A\leq1$ give
\begin{equation}
 \E[P]=\mu,\qquad
 \frac{\operatorname{Var}(P)}{\E[P]^2}
       \leq \frac{\nu}{2^h\mu^2}.                            \label{eq:relvar}
\end{equation}
The numerical sums are truncated after a recorded Poisson Chernoff tail below
$2^{-299}$; that tail is added to the second-moment upper bound.

Chebyshev's inequality yields
\begin{equation}
 \Prb[P<\mu/2]\leq\frac{4\nu}{2^h\mu^2}.                    \label{eq:cheb}
\end{equation}
On the complementary good-table event, the fixed budget
\begin{equation}
                         T=\left\lceil\frac2\mu\right\rceil \label{eq:T}
\end{equation}
succeeds with probability at least $1-e^{-1}$.  Subtracting the signing-cap
abort probability gives a rigorous overall lower bound.  This is a statement
about the realized table, not only expected coverage.

\subsection{Address-prefix sharding}

Fix a $b$-bit prefix before acquisition and retain only addresses with that
prefix.  The retained bucket loads remain independent with the same
$\lambda$.  The success mean becomes $2^{-b}\mu$, the number of stored buckets
falls by $2^b$, and the relative-variance bound in~\eqref{eq:relvar} grows by
$2^b$.  The attack then uses
$T_b=\lceil 2^{b+1}/\mu\rceil$.  This gives a direct time--memory tradeoff
without replaying randomness or predicting acquisition addresses.

\section{Concrete results}

\subsection{All eight submitted sets}

Table~\ref{tab:formal} minimizes $Q+\lceil2/\mu\rceil$.  Each exponent is base
two.  $Q$ is the signing-query Poisson mean and also the number of attacker
$\Hmsg$ recomputations needed to recover the acquired signatures' addresses;
the actual signing count is hard-capped by aborting before $2^{80}$.  $T$ is a
fixed fresh-target $\Hmsg$ budget.  Thus $Q+T$ is attacker public-$\Hmsg$ work,
not a sum of heterogeneous oracle costs.  Every success lower bound exceeds
$0.6321205588$ after subtracting the displayed table tail and the negligible
cap-abort probability.

\begin{table}[ht]
\centering
\small
\caption{Rigorous all-address acquisition and target-search results.}
\label{tab:formal}
\begin{tabular}{lrrrr}
\toprule
set & signer queries $\log_2 Q$ & target $\Hmsg$ $\log_2 T$ &
attacker $\Hmsg$ $\log_2(Q+T)$ & bad-table bound\\
\midrule
160s &74.415&70.593&74.514&$<2^{-62.89}$\\
160f &71.165&66.878&71.237&$<2^{-53.07}$\\
256s &76.670&72.237&76.735&$<2^{-64.23}$\\
256f &72.140&67.248&72.188&$<2^{-54.93}$\\
384s &76.100&71.159&76.146&$<2^{-62.38}$\\
384f &74.430&69.213&74.468&$<2^{-60.30}$\\
512s &77.520&72.367&77.560&$<2^{-62.68}$\\
512f &76.145&70.790&76.180&$<2^{-61.81}$\\
\bottomrule
\end{tabular}
\end{table}

The 160f row's $2^{71.237}$ figure is attacker public-$\Hmsg$ work $Q+T$;
the signing-oracle count is the separate $Q=2^{71.165}$ and the fresh-target
budget is $T=2^{66.878}$.  Table~\ref{tab:physical} reports four separate
physical-resource categories at the same optima.  The table payload
uses a deliberately simple sufficient record: one chain node, one position
byte, and all $a$ authentication nodes for every observed
$(A,i,x)$, plus one complete hypertree suffix per occupied address.  It assumes
no Merkle multiproof compression.  Response and service figures are expected
costs at the listed Poisson means; the donor-table column is expected payload,
exclusive of allocator and database overhead.

\begin{table}[ht]
\centering
\small
\caption{Physical resources at the all-address $Q+T$ optima (base-two
exponents).}
\label{tab:physical}
\scriptsize
\begin{tabular}{lrrrr}
\toprule
set & response bytes / bits & donor bytes & record updates & signer hash calls\\
\midrule
160s &87.745 / 90.745&86.177&78.415&95.802\\
160f &85.410 / 88.410&83.251&75.972&88.767\\
256s &91.312 / 94.312&89.760&81.194&98.659\\
256f &87.542 / 90.542&85.458&77.632&90.939\\
384s &91.989 / 94.989&90.268&81.348&98.451\\
384f &90.647 / 93.647&88.696&80.102&94.938\\
512s &94.101 / 97.101&92.398&83.075&100.852\\
512f &93.105 / 96.105&90.930&82.189&97.388\\
\bottomrule
\end{tabular}
\end{table}

Thus the largest expected exponent in Table~\ref{tab:physical}, in its own
natural unit, is $2^{100.852}$ honest-service hash calls for 512s.  The maxima
for returned data, donor payload, and record updates are separately
$2^{97.101}$ bits, $2^{92.398}$ bytes, and $2^{83.075}$ updates.  The largest
$Q$ and $T$ exponents in Table~\ref{tab:formal} are $77.520$ and $72.367$.
These numbers are neither added nor compared as interchangeable units.

For 160f, a hard worst case after the $2^{80}$ abort contains at most
$2^{94.245}$ returned signature bytes and approximately $2^{97.603}$
PDF-counted honest-signer hash calls.  These worst-case caps are distinct from
the expectations in Table~\ref{tab:physical}.

\subsection{A fully provisioned table below
\texorpdfstring{$2^{80}$}{2-to-the-80} bytes}

The six-bit 160f shard gives the cleanest bounded-memory construction near the
requested resource band.  For one donor record of
\begin{equation}
                        ((a+1)n+1)\ \text{bytes}              \label{eq:record}
\end{equation}
and one address suffix of $(d\,\mathsf{len}+h)n+16$ bytes, a table with every
slot provisioned has deterministic payload
\begin{equation}
 2^{h-6}\left[k2^a((a+1)n+1)+(d\,\mathsf{len}+h)n+16\right]. \label{eq:hardmem}
\end{equation}
Table~\ref{tab:shard} reports both this capacity and the expected occupancy.

\begin{table}[ht]
\centering
\caption{CEDRUS-$\alpha$-160f six-bit-prefix shard.}
\label{tab:shard}
\begin{tabular}{lr}
\toprule
resource & base-two exponent or value\\
\midrule
signing-query mean $Q$ &71.458\\
fixed public-$\Hmsg$ budget $T_6$ &67.160\\
attacker $\Hmsg$ calls $Q+T_6$ &71.530\\
expected full-path donor-table bytes &77.479\\
deterministic fully provisioned payload bytes &79.175\\
returned signature bytes / bits &85.703 / 88.703\\
honest signing-service hash calls &89.060\\
Chebyshev bad-table bound &$<2^{-48.61}$\\
overall success lower bound &$>0.6321205588$\\
\bottomrule
\end{tabular}
\end{table}

The expected returned traffic and signer work do not fall with sharding:
nonretained responses must still be obtained and their $R$ values hashed to
learn their addresses.  The deterministic memory figure is a payload bound,
not a claim that an implementation can allocate $2^{79.175}$ bytes on present
hardware.

\subsection{Alternative physical objectives and fixed-address memory}

Different cost objectives choose different 160f points.  Minimizing the larger
of returned \emph{bits} and public-$\Hmsg$ evaluations gives
$\log_2Q=70.085$, $\log_2T=87.269$, $2^{87.330}$ response bits,
$2^{82.402}$ expected full-path table bytes, and $2^{87.687}$ honest-service
hashes.  Balancing public hashes against signer hashes instead gives
$\log_2Q=70.065$, $\log_2T=87.631$, $2^{87.310}$ response bits, and
$2^{87.667}$ signer hashes.  These are alternative tradeoffs, not resources
paid simultaneously at separate optima.  Each shows why no single
``complexity'' number describes the attack.

A low-memory independent strategy preselects one complete 160f address.  With
$r=554$ retained hits, Paley--Zygmund gives probability at least $0.247483$
that the table covers at least half its mean target mass.  Using
$\lceil2r2^{66}\rceil$ signing queries gives exponent $76.113742$ and a
Chernoff hit-collection failure below $2^{-199.81}$.  A fixed
$2^{72.445408}$ public-target budget yields attacker $\Hmsg$ exponent
$76.222974$ and success at least $0.156439$ minus that failure.  The fully
provisioned one-address payload is only $2^{19.175}$ bytes, but returned
traffic and signer work rise to approximately $2^{90.359}$ bytes and
$2^{93.716}$ hashes.

\section{End-to-end validation}

\subsection{Independent reduced verifier}

The hostile-review implementation uses $h=4$, $a=2$, $k=4$, $W=4$, and
16-byte nodes.  It independently implements domain-separated FORC chains,
complete addresses, leaf-specific authentication paths, an address-fixed FORC
public key, attacker-controlled serialized $R$, ordinary hedged signing, and a
reduced Merkle-authenticated suffix.  SHAKE is an ideal-primitive stand-in at
the reduced node size.

Its fixed run obtained 26 ordinary randomized signatures and forged on a
never-queried message using donors from three distinct messages.  The verifier
accepted the result.  It rejected the same bytes after separately changing the
message, $R$, a FORC authentication node, the hypertree suffix, the chain
direction, or the complete address.

A second, fuller reduced hypertree implementation independently produced both
a one-transcript forward forgery and a coordinate-mixed forgery.  The latter
used two donor transcripts, was accepted on a fresh message, and explicitly
checked that no one acquired transcript covered the target.  Its unadvanced
node, backward-chain, path-tamper, wrong-address, and cross-layer-suffix
controls all rejected.

These are complete reduced experiments, not extrapolated full-parameter runs.
They validate the algebraic splice and all verifier bindings that the proof
uses.

\subsection{Full-parameter FORC-level splice}

A separate deterministic experiment used the exact submitted
$(n,h,a,k,W)$ values, full complete-address widths, Algorithm~16 bit order, and
Algorithms~13--18 address dependencies for all eight parameter sets.  At one
complete address per set it constructed every leaf of every FORC tree, formed
three valid FORC digest openings, selected target coordinates from all three
donors, advanced the exposed nodes, and reused their own authentication paths.
For every set the mixed opening reconstructed the same fixed FORC public key;
no complete donor vector covered the target.  Authentication-path,
unadvanced-node, and wrong-complete-address controls reconstructed a different
key and were rejected.

The experiment was run twice per set: once with the Algorithms~14--18 endpoint
$C(W)$ and once with the prose-shifted endpoint $C(W-1)$.  Every run passed,
confirming computationally that only the order
$\ell_{\rm donor}\leq\ell_{\rm target}$ matters.  SHAKE256 was a deterministic
ideal-primitive stand-in; the experiment validates full submitted FORC
dimensions and the structural splice, not an SM3 implementation and not a
full CEDRUS-$\alpha$ forgery.  It performs neither the predicted acquisition
nor the public target grind.

\subsection{Repaired native 256f conformance}

The submitted source contains three public deviations: 160-bit WOTS message
truncation, eight-bit FORC chain-address truncation, and hash/PRF input
differences~\cite{ngcc-cedrus}.  None is used here.  We repaired a fresh 256f
reference tree to the PDF's full-width address fields, Table~1.2 PRF/hash input
organization, and full constant-sum WOTS message handling.

The native run generated two ordinary hedged signatures on one message and
confirmed that their $R$ values differ.  It checked the published 43,296-byte
signature size; ordinary verification; rejection after changing the message,
$R$, a FORC node, or a WOTS node; rejection of truncated and extended
signatures; rejection under an independent key; a full-width chain-address
test; a fixed Table~1.2 PRF vector; and full-message constant-sum WOTS use.
Every check passed.

The native run is a conformance check, not a full-size forgery.  The full attack
would require the resources in Section~5 and was not executed.

\section{Evaluation boundary and security impact}

The official NGCC Evaluation Criteria, Section~1(2), allow a chosen-message
attacker no more than $2^{80}$ signing queries~\cite{ngcc-criteria}.  The
separate submission requirements, Section~2(2), require a signature key pair
to support at least $2^{64}$ different messages~\cite{ngcc-requirements}.
The latter is a minimum support requirement, not an implicit universal
maximum.

A submission-wide review found no maximum signatures per key, mandatory
key-rotation rule, or $2^{64}$/$2^{80}$ lifetime cap.  The reviewed material
was the 53-page specification, both two-page English and Chinese
basic-information forms, all sixteen reference/optimized implementation
READMEs, all sixteen source trees, and the specification's security-claim
chapter.  The forms contain only identity/contact fields; the byte-identical
READMEs contain generic interface, harness, use, and build notes but no
lifecycle rule.  Chapter~3, Section~3.2 instead says
``let $q_{\rm sig}$ be the number of signing queries overall,'' uses it as an
unconstrained variable in Equations~(3.1)--(3.2), and sums through
$q=q_{\rm sig}$.  The source's 64-bit condition is only an index-width compile
guard, not a usage rule.  The attack samples its query count but
aborts before request $2^{80}$; it therefore respects the evaluation ceiling
in every execution.  For the 160f six-bit shard, the Chernoff upper bound on
that abort event is smaller than $2^{-8.5\cdot10^{24}}$.

The authors' related framework paper makes reduced capacity explicit when it
is intended: its Section~5.5 is titled ``CEDRUS Supporting Less than $2^{64}$
Signatures'' and gives separate CEDRUS+/CEDRUS+FP examples with stated maximum
capacities~\cite{cedrus-framework}.  Its ordinary CEDRUS-$\alpha$ treatment in
Section~5.2 and Appendix~R.3 instead retains a $q_S$-parameterized EU-CMA
reduction without imposing a cap on the NGCC sets.  Those framework parameter
sets are distinct, so this is corroborating documentation practice, not a
claim that Section~5.5 normatively governs the submission.

The distinction is important.  At exactly $Q=2^{64}$, 160f's mean one-target
success is $2^{-161.7365}$.  This method does not break a deployment that
explicitly and enforceably retires each key after $2^{64}$ signatures.  The
submitted construction contains no such rule.  The result is therefore a
parameter and absent-lifetime-policy failure within the stated evaluation
model, rather than a claim that the generic FORC design is unsalvageable.

The impact is an existential fresh-message signature accepted by the ordinary
verifier, not merely a reduced security estimate.  It neither recovers
$\sk.\mathsf{seed}$ nor needs to: the accumulated openings and public grind
directly create a new signature.  The attack is unconditional with respect to
signing state and key choice, while retaining the ideal-primitive assumptions
under which the scheme's concrete analysis is stated.

\section{Scope, limitations, and repair}

The principal limitations are quantitative.
\begin{itemize}[leftmargin=*,itemsep=2pt]
  \item No full-parameter forgery was run.  The formal query counts represent
        astronomically more data and service work than the NREC audit host can
        instantiate.
  \item The probability theorem uses the specification's ideal independent
        $\Hmsg$ model.  It is not an empirical statement about bias in one
        fixed SM3 instantiation.
  \item The simple memory model counts cryptographic payload and explicit
        address metadata.  Real storage systems add indexing and allocator
        overhead.
  \item The attack exceeds $2^{64}$ signing queries.  A lower explicit and
        enforced key lifetime can exclude it.
\end{itemize}

All acquisition queries are ordinary independently randomized chosen-message
signatures.  The construction uses no reset, repeated randomizer, deterministic
mode, related key, weak key, fault, or implementation deviation.  Its material
scope condition is the absence of a lower enforced per-key lifetime: a
deployment that retires the key after $2^{64}$ signatures falls outside the
range established here.

The most direct repair is to state and enforce a per-key signature limit whose
value is included in the concrete security claim.  A $2^{64}$ limit defeats
this route for the submitted rows, although a revised analysis must still
include all public-hash queries and multi-user targets.  Alternatively,
increase the complete address and FORC parameters so the exact quantity
\begin{equation}
  1-\left(1-p_{\rm one}(q_S)\right)^{q_H}                  \label{eq:fullsecurity}
\end{equation}
remains below the target advantage for the claimed $(q_S,q_H)$ range.  The
first and second moments above provide a conservative way to evaluate the
realized multi-address table.

A stateful signer could allocate bottom addresses without reuse, but that
changes the stateless design and requires rollback-safe persistent state.
Merely serializing a public counter does not help if the attacker can vary it
or $R$ during the public target search.  Any forced-pruning or rejection rule
would need a new byte-level specification, distribution analysis, and proof;
none is present in the submitted algorithms.

\section{Prior art and novelty boundary}

The closest prior art is the candidate specification itself.  CEDRUS-$\alpha$
Section~3.2, Equations~(3.1)--(3.2), mixes over address occupancy and gives an
approximate FORC completion probability for one fresh target
\cite{cedrus-spec}.  It already describes the accumulation mechanism.  The
broader CEDRUS framework retains FORC coverage and interleaved target-subset
resilience in its security discussion~\cite{cedrus-framework}.  The SPHINCS+
framework likewise treats degradation across reused few-time instances
\cite{sphincs-framework}.

Abri and Katz's forced-pruning analysis is even closer to the missing concrete
step.  Its Equation~(2) multiplies a one-target success term by the number of
public hash trials $q_{H,2}$ after mixing over bottom-instance occupancy
\cite{abri-katz}.  This is direct prior art for offline public-hash
amplification.  We claim neither that multiplier nor conditioning-safe
coverage counting as new.

At the 3~October~2026 checkpoint, the public NGCC page listed three
CEDRUS-$\alpha$ implementation findings and no parameter/lifetime forgery
\cite{ngcc-cedrus}.  This is a best-effort public-state statement rather than a
proof that no unpublished analysis exists.

The narrow contribution of this report is the candidate-specific conclusion:
exact chain-position coverage and second moments for all eight submitted sets,
a concentration-controlled coordinate-mixing forgery, optimized $q_S/q_H$
tradeoffs with complete physical resource accounting, and the observation that
the absence of a specified per-key lifetime cap places the attack inside the
official NGCC chosen-message query budget.

\section{Conclusion}

Ordinary independently randomized CEDRUS-$\alpha$ signatures can be accumulated
by their complete bottom FORC address.  An attacker who retains the earliest
node and path for each coordinate and leaf can combine openings from different
messages.  Public grinding of serialized $R$ on a fresh message finds a digest
covered by the accumulated table, and one address-matched hypertree suffix
completes a signature accepted by the ordinary verifier.

Exact second moments show that the all-address table concentrates, while a hard
abort keeps signing queries within $2^{80}$.  All eight parameter sets use
fewer than $2^{77.56}$ attacker public-$\Hmsg$ calls $Q+T$, with $Q$ signing
queries reported separately.  The 160f six-bit
shard gives a deterministic donor payload below $2^{80}$ bytes, but still
requires about $2^{85.70}$ response bytes and $2^{89.06}$ honest-service hash
calls.  At the all-address optima, 512s reaches the largest expected natural
unit exponent, $2^{100.852}$ honest-signer hashes.  The attack is therefore a
concrete parameter/lifetime failure in the evaluation model, not a
machine-practical computation.

CEDRUS-$\alpha$ already identifies FORC accumulation, and Abri--Katz already
identify the public-query multiplier.  The repair is to make the claimed key
lifetime and public-hash budget explicit and choose or enforce parameters for
their joint effect.

\section*{Acknowledgements}

The computations were performed on the Norwegian Research and Education Cloud
(NREC), using resources provided by the University of Bergen and the University
of Oslo.

\section*{Artifact availability}

The accompanying reproducer contains both reduced verifiers, the exact
all-parameter probability and resource calculation, the all-eight-set
full-parameter FORC splice, a bundled repaired 256f native source tree, the
submission-wide lifetime/counter audit, negative controls, expected
machine-readable outputs, and a one-command replay.  It requires no network
access.  The artifact verifies its immutable inputs before running and confines
generated files to disposable build and work directories.

\begin{thebibliography}{99}
\small

\bibitem{cedrus-spec}
Siwei Sun, Zhen Qin, Fangyu Zheng, and Zhiyu Zhang,
\emph{CEDRUS$\alpha$: A Stateless Hash-based Signature Scheme},
NGCC submission specification, 29 June 2026.

\bibitem{cedrus-framework}
Zhen Qin, Siwei Sun, and Fangyu Zheng,
\emph{Extending the SPHINCS+ Framework: Varying the Tree Heights and Chain
Lengths}, IACR Cryptology ePrint Archive, Report 2025/2236, revised 6 September
2026,
\url{https://eprint.iacr.org/2025/2236}.

\bibitem{abri-katz}
Mehdi Abri and Jonathan Katz,
\emph{Shorter Hash-Based Signatures Using Forced Pruning},
IACR Cryptology ePrint Archive, Report 2025/2069, revised 2026; CRYPTO 2026,
\url{https://eprint.iacr.org/2025/2069}.

\bibitem{sphincs-framework}
Daniel J. Bernstein, Andreas H\"ulsing, Stefan K\"olbl, Ruben Niederhagen,
Joost Rijneveld, and Peter Schwabe,
\emph{The SPHINCS+ Signature Framework},
IACR Cryptology ePrint Archive, Report 2019/1086, 2019,
\url{https://eprint.iacr.org/2019/1086}.

\bibitem{ngcc-criteria}
National Cryptography Administration of China,
\emph{Evaluation Criteria for Candidate Algorithms of the Next-Generation
Public-Key Cryptography Competition}, Section~1(2), 2025,
\url{https://www.niccs.org.cn/niccs/Notice/1975896137741635584/tT7TSQiz.pdf}.

\bibitem{ngcc-requirements}
National Cryptography Administration of China,
\emph{Submission Requirements for the Next-Generation Public-Key
Cryptography Competition}, Section~2(2), 2025,
\url{https://www.niccs.org.cn/niccs/Notice/lDop1mav.pdf}.

\bibitem{ngcc-cedrus}
NGCC Security Reports,
\emph{CEDRUS-$\alpha$ (sign-04)},
\url{https://ngcc.dev/reports/sign-04.html}, accessed 3 October 2026.

\end{thebibliography}

\end{document}
